\documentclass[12pt]{article}
\usepackage[utf8]{inputenc}
\usepackage[T1]{fontenc}
\usepackage[margin=1in]{geometry}
\usepackage{setspace}
\usepackage{xcolor}
\usepackage[colorlinks=true, urlcolor=blue!60!black, linkcolor=blue!60!black]{hyperref}
\setlength{\parindent}{2em}
\setlength{\parskip}{1.25em}
\renewcommand{\baselinestretch}{1.0}
\begin{document}
\begin{center}
{\LARGE \textit{Myers v. United States}, 272 U.S. 52 (1926)\par}
\vspace{0.8em}
{\large Case Brief --- Four-Page Edition\par}
\end{center}
\vspace{1em}

\section{Memory Jogger}

When President Wilson fired Portland postmaster Frank Myers without the Senate consent an 1876 statute required, Chief Justice Taft's 6--3 opinion held that Article II gives the President exclusive power to remove executive officers he appoints with Senate consent, so Congress could not require the Senate's consent to removal (\href{https://supreme.justia.com/cases/federal/us/272/52/}{272 U.S. 52}).

\section{Detailed Case Facts}

On July 21, 1917, President Wilson, with the Senate's advice and consent, appointed Frank S.\ Myers to a four-year term as first-class Postmaster in Portland, Oregon. An 1876 statute provided that such postmasters ``shall be appointed and may be removed by the President by and with the advice and consent of the Senate.''

In January 1920 the Post Office Department demanded Myers's resignation, and he refused. On February 2, 1920, the Postmaster General, acting at Wilson's direction, removed him without seeking the Senate's consent. Myers petitioned the President and the Senate committee for a hearing and protested until his term expired on July 21, 1921, taking no other work.

\section{Procedural History}

Myers sued in the Court of Claims for \$8,838.71 in lost salary; after his death his administratrix continued the case. The Court of Claims ruled for the United States on laches. On direct appeal, the Supreme Court rejected laches, ordered reargument, and invited Senator George Wharton Pepper to defend the statute as \textit{amicus curiae}. Argued in 1924 and reargued in 1925, the case was decided on October 25, 1926.

\section{Judicial Votes}

Affirmed on constitutional grounds, \textbf{6--3}.
\begin{itemize}
\item \textbf{Majority:} Chief Justice Taft, joined by Justices Van Devanter, Sutherland, Butler, Sanford, and Stone.
\item \textbf{Concurring:} None.
\item \textbf{Dissenting:} Justice Holmes; Justice McReynolds; Justice Brandeis (each separately).
\end{itemize}

\section{Holding}

Article II grants the President exclusive power to remove executive officers he appoints with the Senate's advice and consent, and Congress may not condition that power on Senate concurrence. The 1876 postmaster provision, like the Tenure of Office Act of 1867, is unconstitutional. The judgment for the United States was affirmed.

\section{Analysis of Opinions}

\textbf{Majority (Taft).} Drawing on his own presidency, Taft contrasted Article I's ``legislative Powers herein granted'' with Article II's general vesting of ``executive Power.'' Following Hamilton's \textit{Pacificus}, he read Article II as a general grant subject to express exceptions. Because the President must execute the laws through subordinates and must take care that they be faithfully executed, he needs ``general administrative control of those executing the laws,'' including power to remove them. Removal is incident to appointment, but the Senate's advice-and-consent role is a narrow exception that cannot be extended to removals: blocking a nominee only requires another choice, while blocking a removal forces the President to keep a disloyal subordinate.

Taft relied heavily on the Decision of 1789, in which the First Congress, led by Madison, framed the foreign-affairs statute to imply a constitutional removal power, and on 74 years of acquiescence. He treated the 1867 and 1876 statutes as partisan aberrations. Congress may regulate removal of inferior officers only if it vests their appointment in someone other than the President, and may never draw removal power to itself. Even officers with quasi-judicial duties may be removed for how they exercised discretion.

\textbf{Holmes (dissenting).} Taft's deductions were ``spider's webs inadequate to control the dominant facts''; Congress creates the office and may condition removal, and the take care duty does not go beyond the laws.

\textbf{McReynolds (dissenting).} No express clause grants an illimitable removal power; ouster is executive, but prescribing tenure conditions is legislative.

\textbf{Brandeis (dissenting).} Reviewing 137 years of practice, Brandeis showed Congress had often restricted removal of inferior officers, and argued that the separation of powers aims to prevent arbitrary power, not to promote efficiency.

\textbf{Later developments.} \textit{Humphrey's Executor} (1935) confined \textit{Myers} to purely executive officers; \textit{Seila Law} revived it; and \textit{Trump v.\ Slaughter} (2026) overruled \textit{Humphrey's Executor}.

\section{Examples: Future Applications}

\begin{enumerate}
\item \textbf{Same-Side (Invalid): A Federal AI Commission.} Commissioners who issue binding rules, adjudicate, and seek penalties are protected from removal except for cause. After \textit{Slaughter}, such executive officers must be removable at will.
\item \textbf{Same-Side (Invalid): Congressional Removal.} A Comptroller who certifies budget cuts is removable only by joint resolution. Congress may not draw removal power to itself, as \textit{Bowsher} held.
\item \textbf{Opposite-Side (Valid): A Department-Appointed Inferior Officer.} A naval officer appointed by the Secretary of the Navy may be dismissed in peacetime only by court-martial. \textit{Myers} itself distinguished this \textit{Perkins} situation.
\item \textbf{Opposite-Side (Not Triggered): A Civil-Service Examiner.} An ordinary employee is not an officer under \textit{Buckley v.\ Valeo}, and \textit{Free Enterprise Fund} left the civil service untouched.
\item \textbf{Fence-Sitter: The Commission's Chief Inspector.} An inferior officer appointed by the Chair has single-layer for-cause protection. \textit{Perkins} and \textit{Morrison} support it once the Chair is removable at will, but a maximalist reading of \textit{Myers} could reach officers wielding core enforcement power.
\end{enumerate}

\section{Critique}

\textbf{Contested history.} Modern legal historians, cited by Justice Sotomayor in her \textit{Slaughter} dissent, argue that Taft overstated the First Congress's consensus. The 1789 House was split into three or four factions, and its rejection of Senate participation in removal did not necessarily settle whether Congress could impose for-cause tenure limits by statute.

\textbf{Concentrated power.} As Justice Gorsuch noted in \textit{Slaughter}, combining \textit{Myers}'s plenary removal power with broad modern delegations of rulemaking and adjudication concentrates unprecedented power in one executive, a tension that neither Taft in 1926 nor the \textit{Slaughter} majority fully resolved.

\section{Key Quotations}

\begin{quote}
``The doctrine of the separation of powers was adopted by the convention of 1787, not to promote efficiency but to preclude the exercise of arbitrary power. The purpose was not to avoid friction, but, by means of the inevitable friction incident to the distribution of the governmental powers among three departments, to save the people from autocracy.''
\end{quote}
\noindent --- Justice Brandeis, dissenting, 272 U.S.\ at 293.

\begin{quote}
``The vesting of the executive power in the President was essentially a grant of the power to execute the laws. But the President alone and unaided could not execute the laws. He must execute them by the assistance of subordinates. ... As he is charged specifically to take care that they be faithfully executed, the reasonable implication, even in the absence of express words, was that as part of his executive power he should select those who were to act for him under his direction in the execution of the laws. The further implication must be, in the absence of any express limitation respecting removals, that as his selection of administrative officers is essential to the execution of the laws by him, so must be his power of removing those for whom he can not continue to be responsible.''
\end{quote}
\noindent --- Chief Justice Taft, Opinion of the Court, 272 U.S.\ at 117.

\begin{quote}
``Then there may be duties of a quasi judicial character imposed on executive officers and members of executive tribunals whose decisions after hearing affect interests of individuals, the discharge of which the President can not in a particular case properly influence or control. But even in such a case he may consider the decision after its rendition as a reason for removing the officer, on the ground that the discretion regularly entrusted to that officer by statute has not been on the whole intelligently or wisely exercised. Otherwise he does not discharge his own constitutional duty of seeing that the laws be faithfully executed.''
\end{quote}
\noindent --- Chief Justice Taft, Opinion of the Court, 272 U.S.\ at 135.

\end{document}
